\chapter{Related Work}
\label{chap:related_work}

A moral label attached to an image appears to name something in the scene, but
it compresses two judgements: whether the scene is morally significant and
which concern best explains that significance. This chapter follows that
problem from moral categories, through visual ambiguity and existing benchmark
designs, to the comparison made in this dissertation.


\section{Moral Categories as a Comparative Vocabulary}
\label{sec:moral_foundations}

Moral judgement can concern whether an act is wrong, a norm has been violated,
or an agent deserves blame \cite{malle2021moral}. An image can also be morally
meaningful without showing wrongdoing, for example when it depicts care,
commitment, or respect.

Moral Foundations Theory (MFT) organises such concerns into five widely used
virtue--vice pairs: Care/Harm, Fairness/Cheating, Loyalty/Betrayal,
Authority/Subversion, and Sanctity/Degradation
\cite{graham2009liberals,graham2011mapping}. They describe different reasons a
situation might matter morally, not positions on one scale from good to bad.

This vocabulary supports comparison, but does not make its labels objective
properties of images. The MFQ-2, for example, is a validated six-foundation
self-report measure, not an image-classification instrument
\cite{zakharin2023mfq2}. Relationships among foundations also vary across
cultures \cite{atari2023weird}. A picture may therefore support several
readings, or too little evidence for any of them.

The present study adapts the five pairs into ten image categories and adds
\textit{None of these / Not clearly moral in nature}. MFT is thus a shared
comparative language, not a complete or culturally neutral account of visual
morality. The model must still prioritise one label when several are plausible.


\section{Why Images Create a Two-Stage Moral Problem}
\label{sec:visual_moral_inference}

Many computational studies of morality begin with text. ETHICS, Moral Stories,
and Delphi can state actions, intentions, norms, and consequences directly
\cite{hendrycks2021aligning,emelin2021moral,jiang2025delphi}. A single image
often leaves that context unstated.

An image may show an outcome without establishing what preceded it, what the
people intended, how they are related, or what follows. Visual commonsense
research distinguishes recognising visible objects and actions from inferring
goals and social relations \cite{zellers2019recognition}. The distinction
becomes morally important when judgement depends on intention or responsibility
\cite{malle2021moral}. Visible distress, for example, does not by itself reveal
whether a scene concerns Harm, Care, Betrayal, or none of these.

\textbf{A model must first decide whether the evidence crosses a
moral-relevance boundary, and only then which moral frame should dominate.}
Two models can agree at the first stage but disagree at the second because they
give different weight to expression, posture, context, or an inferred story.

The Socio-Moral Image Database (SMID) preserves this complexity through
continuous ratings. Its 2,941 photographs received more than 820,000
judgements from 2,716 participants on affect, moral wrongness, and foundation
relevance \cite{crone2018smid}. Using SMID, Zhu et al. found that visual
representations added predictive information not fully captured by generated
captions \cite{zhu2025visual}. In that setting, a description was not a
complete substitute for the image.

These studies use richer targets than the forced choice examined here. That
choice removes graded relevance and competing labels, but also reveals where a
model draws the moral boundary and which reading it prioritises.


\section{How Multimodal Studies Handle Ambiguity}
\label{sec:multimodal_moral_evaluation}

Recent multimodal studies respond to ambiguous moral evidence through
different design choices. MM-MoralBench combines synthetic visual contexts
with dialogue and evaluates judgement, classification, and response tasks
across six foundations \cite{yan2026mmmoralbench}. The controlled scenarios
and dialogue provide information that an unaccompanied image may omit.
MORALISE instead uses 2,481 expert-verified real-world image--text pairs and a
13-topic taxonomy grounded in Turiel's Domain Theory
\cite{lin2026moralise}. It separates detecting a moral violation from
identifying the norm involved. This distinction closely resembles the
two-stage problem described above, although MORALISE begins with annotated
moral-violation material and uses a different taxonomy.

Other work preserves ambiguity by changing the response rather than adding
context. Rezapour et al. asked people and language models to construct moral
stories from shared visual cues, finding differences in both moral emphasis
and narrative form \cite{rezapour2025tales}. Their open-ended format allows a
scene to develop into several possible stories. MM-SCALE takes a different
route, using five-point acceptability ratings, grounded reasoning labels, and
listwise preferences \cite{park2026mmscale}. It therefore offers a richer
target than a binary or single-label judgement.

Table~\ref{tab:literature_design_comparison} brings the central design contrast
together. These resources differ not only in scale, but in what uncertainty
their recorded targets preserve. The present study has no corpus-wide,
independently validated reference structure. Its purpose is consequently not
to reproduce benchmark accuracy, but to examine what happens when deployed
models face the same restricted label space on a mixed set that also permits
\textit{None}.

\begin{table}[!htbp]
    \centering
    \caption{How closely related visual moral studies represent context and
    judgement.}
    \label{tab:literature_design_comparison}
    \footnotesize
    \begin{spacing}{1.0}
    \setlength{\tabcolsep}{4pt}
    \renewcommand{\arraystretch}{1.12}
    \begin{tabularx}{\textwidth}{@{}L{0.18\textwidth}YYY@{}}
        \toprule
        Resource & Visual material & Recorded target & Main design choice \\
        \midrule
        SMID \cite{crone2018smid} & 2,941 photographs & Graded human ratings
        of wrongness, affect, and foundation relevance & Retains degrees of
        relevance and observer variation \\
        MM-MoralBench \cite{yan2026mmmoralbench} & Synthetic visual contexts
        with dialogue & Foundation-based judgement, classification, and
        response tasks & Adds controlled moral context \\
        MORALISE \cite{lin2026moralise} & 2,481 curated image--text pairs &
        Moral judgement and 13-topic norm attribution & Separates detecting a
        violation from identifying its norm \\
        MM-SCALE \cite{park2026mmscale} & Image--scenario pairs & Five-point
        acceptability and listwise preferences & Replaces a single target with
        graded, grounded comparisons \\
        Present study & 127 mixed-source images & One model label and
        confidence; participant, repeat, and edit subsets & Audits a shared
        forced-choice task without treating one label as correct \\
        \bottomrule
    \end{tabularx}
    \end{spacing}
\end{table}

\FloatBarrier

Multimodal safety benchmarks likewise show that images matter when models
handle harmful inputs \cite{liu2024mmsafetybench,ying2024safebench}. Their
outcome is safe response behaviour, however, not the moral categorisation of
an ambiguous scene; the two alignment questions are related but distinct.


\section{Agreement, Human Variation, and Reported Confidence}
\label{sec:model_uncertainty}

Once an image becomes a label, the next question is what agreement with it
means. Text-based studies show that elicitation can shape moral outputs.
Abdulhai et al. examined moral-foundation profiles across conversational and
targeted prompt contexts \cite{abdulhai2024foundations}; Liu et al. compared
moral choices, ranked principles, and choice firmness across models and Chinese
university students \cite{liu2024pluralistic}. These studies treat an output
as a response to an instrument, not direct access to fixed values. Alternative
prompts should likewise be robustness checks, not a search for a model's
``true'' morality.

The same caution applies to confidence. In text-based knowledge and reasoning
tasks, its elicitation affects how well it corresponds to performance
\cite{xiong2024uncertainty}. A requested confidence score is therefore not
automatically a calibrated probability. It may show how strongly a model
presents its answer without showing how widely others share that answer.

Human agreement is not a simple substitute for model agreement. The importance
assigned to moral foundations differs across people, groups, and cultural
contexts \cite{graham2009liberals,graham2011mapping,qian2024cultural,
atari2023weird}. VLM image understanding also varies across cultural contexts
\cite{yadav2025beyond,ananthram2024perspective}. This literature does not
explain a particular case here, but cautions against treating one participant
sample as a universal human standard.

Subjective-annotation research further shows that majority voting can erase
systematic differences between annotators \cite{davani2022dealing}. Response
distributions may therefore be evidence rather than noise
\cite{uma2021disagreement}. A plurality label cannot show whether most people
agreed or several readings received substantial support.

\textbf{Agreement consequently has several meanings.} Models can agree with
one another, a model can reproduce itself, or participants can include its
label. Confidence is another quantity again. None validates the others.


\section{Visual-Cue Sensitivity and the Present Study}
\label{sec:research_positioning}

Aggregate agreement does not explain why a particular label was selected.
Comparing an image with an edited version can make possible cue dependence more
visible. Liu et al. tested VLM moral judgements under textual and visual
perturbations and found that robustness varied by perturbation, moral domain,
and model \cite{liu2026moralbackbone}. Their work shows the value of evaluating
whether a judgement survives a changed input. It also sets a boundary on what
small edit studies can establish: if an edit changes several visible
properties at once, a changed label can motivate a hypothesis about relevant
cues but cannot isolate a causal mechanism.

Taken together, these studies leave open the combined question examined here:
how several deployed LVLM configurations use the same single-label MFT
vocabulary on a fixed corpus containing both morally suggestive and ambiguous
images, and how their outputs differ across label resolution, repeated runs,
participant multi-label responses, and matched image edits.

The four research questions follow that sequence. RQ1 asks where each model
draws the boundary between a moral category and \textit{None}, and which frames
appear in its retained outputs. RQ2 asks whether apparent agreement survives
when the labels are examined at finer resolution and when classifications are
repeated. RQ3 keeps the distribution of participant selections visible, so a
model's correspondence with a plurality can be distinguished from broad human
agreement. RQ4 then uses selected edits to identify cases in which changing the
visible evidence is followed by a different label or confidence score. A
supplementary prompt experiment tests whether the participant-correspondence
pattern is robust to alternative evidence-checking instructions; it is not a
fifth research question.

Measurement and benchmark research stresses that an evaluation result inherits
the assumptions and scope of the task that produced it
\cite{jacobs2021measurement,raji2021benchmark}. This dissertation therefore
does not turn model consensus, participant plurality, researcher indexing, or
confidence into an accuracy target. \textbf{Its contribution is to show what a
simple one-label comparison can reveal---the moral boundary, the preferred
frame, stability, participant correspondence, and sensitivity to edited
evidence---and what that comparison would hide if these were reduced to one
score.}
